../cictikz/src/cictikz/data/tex/cictikz_preamble.tex

% A quadrature continuous-time sigma-delta modulator, drawn the way the
% text asks you to read it rather than the way it is built.
%
% Inspired by the modulator in "A 56 mW Continuous-Time Quadrature
% Cascaded Sigma-Delta Modulator With 77 dB DR in a Near Zero-IF 20 MHz
% Band". The original is fully differential, cascaded, and carries every
% tuning resistor; the surrounding text tells the reader to set the
% cross-coupling to zero, follow one path, and ignore the second stage,
% so that is what this draws, with the cross-coupling and the second
% stage shown as the two things you add back afterwards.

\begin{tikzpicture}[thick]
  %%%%%%%%%%%%%%%%%%%%%%%%%%%%%%%%%%%%%%%%%%%%%%%%%%%%%%%%%%%%%%%%%%%%%%
%% Shared pieces for the l04_afe signal flow graphs, l4_first_order and
%% l4_biquad.  The biquad is meant to read as the first-order graph with a
%% second integrator added, so the summing node and the 1/s block have to be
%% the same size and shape in both.  They are the same size in the original
%% artwork too.
%%
%% Names are prefixed and letters only, and computed coordinates are braced:
%% TikZ scans a coordinate for its closing parenthesis, so a '(' inside an
%% expression ends it early.
%%%%%%%%%%%%%%%%%%%%%%%%%%%%%%%%%%%%%%%%%%%%%%%%%%%%%%%%%%%%%%%%%%%%%%

\newcommand{\sfgR}{0.25}    % summing node radius
\newcommand{\sfgBw}{1.0}    % half width of the 1/s block
\newcommand{\sfgBh}{0.65}   % half height of the 1/s block

% A summing node at (#1,#2).
\newcommand{\sfgSum}[2]{%
  \draw (#1,#2) circle (\sfgR);
  \draw ({#1-0.17},#2) -- ({#1+0.17},#2);
  \draw (#1,{#2-0.17}) -- (#1,{#2+0.17});
}

% A block centred at (#1,#2) carrying #3.
\newcommand{\sfgBox}[3]{%
  \draw ({#1-\sfgBw},{#2-\sfgBh}) rectangle ({#1+\sfgBw},{#2+\sfgBh});
  \node at (#1,#2) {#3};
}

\newcommand{\sfgDot}[2]{\fill (#1,#2) circle (0.07);}


  \def\yi{1.9}
  \def\yq{-1.9}

  \foreach \y/\tag/\sg in {\yi/i/1, \yq/q/-1} {
    \draw[->] (-3.6,\y) -- (-2.47,\y);
    \node[anchor=east, scale=0.9] at (-3.7,\y) {$V_{\tag}$};
    \node[scale=0.85, anchor=south] at (-3.1,\y+0.1) {$R_1$};

    \sfgSum{-2.2}{\y}
    \sfgBox{0.4}{\y}{$\displaystyle\frac{1}{s}$}
    \draw[->] (-1.95,\y) -- (-0.6,\y);

    \sfgBox{3.6}{\y}{$H_2(s)$}
    \draw[->] (1.4,\y) -- (2.6,\y);

    \draw[->] (4.6,\y) -- (5.6,\y);
    \sfgTall{5.6}{\y-0.7}{6.9}{\y+0.7}{ADC};
    \draw[->] (6.9,\y) -- (8.3,\y);
    \node[anchor=west, scale=0.9] at (8.4,\y) {$Y_{\tag}$};
    %- not at 7.6: that is where the feedback taps off, so the label sat
    %- on top of the vertical wire
    \node[scale=0.8, anchor=south] at (7.2,\y+0.1) {4b};

    %- the feedback that makes it a sigma-delta at all
    \draw (7.6,\y) -- (7.6,{\y+\sg*1.5});
    \fill (7.6,\y) circle (0.075);
    \sfgTall{2.9}{\y+\sg*1.5-0.4}{4.3}{\y+\sg*1.5+0.4}{DAC};
    \draw[->] (7.6,{\y+\sg*1.5}) -- (4.3,{\y+\sg*1.5});
    \draw (2.9,{\y+\sg*1.5}) -- (-2.2,{\y+\sg*1.5});
    \draw[->] (-2.2,{\y+\sg*1.5}) -- (-2.2,{\y+\sg*0.3});
  }

  %- The cross-coupling, tapped from each integrator output and taken to
  %- the other path's summing node.
  %-
  %- The two arms cannot both drop straight down the middle: each summing
  %- node already has its input on the left and its DAC on the outer
  %- side, which leaves only the inner side, and both arms want it. Run
  %- at x = -2.35 and -2.05 they were three tenths apart and read as one
  %- thick wire. So the I to Q arm takes the inner approach, and the Q to
  %- I arm goes out past the input on the far left and comes back in
  %- underneath it.
  \draw[red, thick] (1.4,\yi) -- (1.4,0.85) -- (-2.2,0.85);
  \draw[red, thick, ->] (-2.2,0.85) -- (-2.2,{\yq+0.24});
  \fill[red] (1.4,\yi) circle (0.075);

  \draw[red, thick] (1.4,\yq) -- (1.4,-0.85) -- (-3.4,-0.85)
    -- (-3.4,1.5);
  \draw[red, thick, ->] (-3.4,1.5) -- (-2.36,{\yi-0.21});
  \fill[red] (1.4,\yq) circle (0.075);

  \node[red, scale=0.85, anchor=west] at (1.7,0.0)
    {$R_{fb}$: the cross-coupling};

  \node[anchor=north, align=center, scale=0.9] at (2.6,-4.4)
    {Set $R_{fb}$ to infinity and the two halves are ordinary real\\
     sigma-delta modulators. Connect them, and the noise transfer\\
     function stops being symmetric about zero frequency, which is\\
     what a near zero-IF receiver needs.};

\end{tikzpicture}
\end{document}
